\plainchapter{Exam Pattern \& the 250+ Strategy}{intro}

\needspace{95pt}

\hypertarget{intro--1-exam-pattern-at-a-glance}{%
\section{Exam Pattern at a Glance}\label{intro--1-exam-pattern-at-a-glance}}

\diagram{figures/intro-00}

\begingroup\small

\begin{longtable}[]{@{}
  >{\raggedright\arraybackslash}p{(\columnwidth - 4\tabcolsep) * \real{0.1678}}
  >{\raggedright\arraybackslash}p{(\columnwidth - 4\tabcolsep) * \real{0.3523}}
  >{\raggedright\arraybackslash}p{(\columnwidth - 4\tabcolsep) * \real{0.2023}}@{}}
\toprule\noalign{}
\begin{minipage}[b]{\linewidth}\raggedright
\textbf{Item}
\end{minipage} & \begin{minipage}[b]{\linewidth}\raggedright
\textbf{Paper 1}
\end{minipage} & \begin{minipage}[b]{\linewidth}\raggedright
\textbf{Paper 2}
\end{minipage} \\
\midrule\noalign{}
\endhead
\bottomrule\noalign{}
\endlastfoot
Questions & 50 (all compulsory) & 100 (all compulsory) \\
Marks & 100 & 200 \\
Negative marking & None & None \\
Nature & Teaching/\allowbreak{}Research aptitude, reasoning & Core CS (10 units) \\
\end{longtable}

\endgroup

\textbf{No negative marking → never leave a question blank.}

\needspace{125pt}

\hypertarget{intro--2-what-250-really-means}{%
\section{What 250+ Really Means}\label{intro--2-what-250-really-means}}

250 / 300 = \textbf{125 correct out of 150} (≈ 83\% accuracy).

\diagram{figures/intro-01}

\begingroup\small

\begin{longtable}[]{@{}
  >{\raggedright\arraybackslash}p{(\columnwidth - 8\tabcolsep) * \real{0.0920}}
  >{\raggedright\arraybackslash}p{(\columnwidth - 8\tabcolsep) * \real{0.1215}}
  >{\raggedright\arraybackslash}p{(\columnwidth - 8\tabcolsep) * \real{0.1672}}
  >{\raggedright\arraybackslash}p{(\columnwidth - 8\tabcolsep) * \real{0.0794}}
  >{\raggedright\arraybackslash}p{(\columnwidth - 8\tabcolsep) * \real{0.2080}}@{}}
\toprule\noalign{}
\begin{minipage}[b]{\linewidth}\raggedright
\textbf{Paper}
\end{minipage} & \begin{minipage}[b]{\linewidth}\raggedright
\textbf{Questions}
\end{minipage} & \begin{minipage}[b]{\linewidth}\raggedright
\textbf{Target correct}
\end{minipage} & \begin{minipage}[b]{\linewidth}\raggedright
\textbf{Marks}
\end{minipage} & \begin{minipage}[b]{\linewidth}\raggedright
\textbf{Allowed mistakes}
\end{minipage} \\
\midrule\noalign{}
\endhead
\bottomrule\noalign{}
\endlastfoot
Paper 1 & 50 & 42+ & 84+ & 8 \\
Paper 2 & 100 & 84+ & 168+ & 16 \\
\textbf{Total} & \textbf{150} & \textbf{126} & \textbf{252} & \textbf{24} \\
\end{longtable}

\endgroup

\needspace{130pt}

\hypertarget{intro--where-marks-come-from-approx-paper-2-weightage--10-qs-per-unit}{%
\subsection{Where marks come from (approx. Paper 2 weightage --- \textasciitilde10 Qs per unit)}\label{intro--where-marks-come-from-approx-paper-2-weightage--10-qs-per-unit}}

\begingroup\small

\begin{longtable}[]{@{}
  >{\raggedright\arraybackslash}p{(\columnwidth - 6\tabcolsep) * \real{0.3387}}
  >{\raggedright\arraybackslash}p{(\columnwidth - 6\tabcolsep) * \real{0.1445}}
  >{\raggedright\arraybackslash}p{(\columnwidth - 6\tabcolsep) * \real{0.1199}}
  >{\raggedright\arraybackslash}p{(\columnwidth - 6\tabcolsep) * \real{0.2226}}@{}}
\toprule\noalign{}
\begin{minipage}[b]{\linewidth}\raggedright
\textbf{Unit}
\end{minipage} & \begin{minipage}[b]{\linewidth}\raggedright
\textbf{Expected Qs}
\end{minipage} & \begin{minipage}[b]{\linewidth}\raggedright
\textbf{Difficulty}
\end{minipage} & \begin{minipage}[b]{\linewidth}\raggedright
\textbf{Priority for 250+}
\end{minipage} \\
\midrule\noalign{}
\endhead
\bottomrule\noalign{}
\endlastfoot
Discrete Structures \& Optimization & 10--12 & Medium & ★★★ \\
Computer System Architecture & 8--10 & Medium & ★★★ \\
Programming Languages \& Graphics & 8--10 & Easy-Med & ★★ \\
DBMS & 10--12 & Easy-Med & ★★★ (scoring) \\
System Software \& OS & 10--12 & Medium & ★★★ (scoring) \\
Software Engineering & 8--10 & Easy & ★★★ (scoring, theory) \\
Data Structures \& Algorithms & 10--12 & Medium & ★★★ \\
TOC \& Compilers & 10--12 & Hard & ★★★ (differentiator) \\
Networks & 10--12 & Medium & ★★★ \\
Artificial Intelligence & 8--10 & Medium & ★★ \\
\end{longtable}

\endgroup

Paper 1: roughly 5 questions per unit; \textbf{Data Interpretation, Reasoning, Comprehension} are "sure marks" (≈ 15 Qs) --- aim for 100\% there.

\needspace{161pt}

\hypertarget{intro--5-about-the-pyqs}{%
\section{About the PYQs}\label{intro--5-about-the-pyqs}}

The questions under \textbf{"Previous Year Questions (PYQ pattern)"} are reconstructed from the recurring
question types and concepts asked in NTA UGC NET (2012--2025) and its predecessor CBSE NET papers. Numbers and
wording may differ from the original papers, and year tags are deliberately left out where they could not be verified.
Always pair these notes with the \textbf{official NTA question papers \& answer keys} (ugcnet.nta.ac.in) for the
last 5--6 cycles --- solve them fully timed at least twice.

\needspace{95pt}

\hypertarget{intro--6-golden-rules-for-250}{%
\section{Golden Rules for 250+}\label{intro--6-golden-rules-for-250}}

\begin{enumerate}
\def\labelenumi{\arabic{enumi}.}
\tightlist
\item
  \textbf{Paper 2 decides your score} --- 2/3 of total marks. Spend \textasciitilde70\% of your time on it.
\item
  \textbf{Theory units are free marks} (SE, DBMS, OS, Networks, AI): learn definitions exactly.
\item
  \textbf{Numerical units are differentiators} (Discrete, COA, TOC, Algorithms): practise daily.
\item
  \textbf{PYQs repeat} --- concepts recur cycle after cycle. Solve every PYQ in these files twice.
\item
  \textbf{Mocks}: at least 25 full-length mocks in the last 6 weeks; analyse every wrong answer.
\item
  \textbf{Attempt all 150} --- no negative marking.
\end{enumerate}
